\chapter{算法验证与性能评估} 

本章旨在对前文提出的非全维位置量测条件下的分布式状态观测器与鲁棒协同控制策略进行验证与性能评估。本章首先通过 MATLAB 数值仿真平台，针对四旋翼无人机编队的系统可观测性、状态估计收敛性以及机间避碰与通信连通性等关键理论指标进行严谨验证；随后，利用更贴近真实物理环境的 Gazebo 三维动力学仿真框架，为算法从理论走向实际工程部署提供进一步的验证思路与方案。

\section{MATLAB仿真实验与结果分析}
为验证所提编队控制协议的有效性，本节在 MATLAB 环境中搭建了基于欧拉数值积分法的多无人机系统仿真平台。

\subsection{仿真参数设置与拓扑配置}
仿真场景设定为三维空间中由$N=4$架无人机组成的微型集群系统，其动力学行为由公式(\ref{一阶动力学})所描述。同时，假设每架无人机所配备的底层控制器或执行器能够快速且精确地跟踪期望速度指令$\mathbf{v}_i^*$，从而支撑上层编队控制器的正常运行。

首先，验证编队局部收敛的具体配置与章节\ref{异构传感器配置实例分析}一致。无人机集群的底层通信与传感拓扑为全连通图，共包含$M=6$条双向通信链路。为了满足第3章提出的系统弱可观测性要求，对4架无人机进行了异构传感器配置，即部分无人机仅能量测特定轴向的绝对位置，而另一部分仅具备相对测距能力。系统的期望编队构型被设定为一个边长为$\sqrt{2}$m的正四面体。编队$\mathbf{p}^*$在全局坐标系下的位置分别为$\mathbf{p}_1^*=[1;1;1]$，$\mathbf{p}_2^*=[2;1;2]$，$\mathbf{p}_3^*=[1;2;2]$，$\mathbf{p}_4^*=[2;2;1]$，单位均为m。测距传感器网络与通信网络相同，均建模为无向图，由边$\mathcal{E}=\{(i,j) \mid i,j\in \mathcal{I}_{4}\}$表征。另外，仿真设置中还包含了一些增益参数，详情见表\ref{无人机集群系统刚性编队仿真基础参数表}。根据距离约束可以计算出无人机初始位置的邻域参数$\rho_p=0.29$，基于此可以对无人机的初始位置$\mathbf{p}_i(0)$进行配置，初始位置估计$\hat{\mathbf{p}}_i(0)$根据初始位置随即给出。此外，为验证系统在初始偏差下的鲁棒收敛能力，所有无人机的初始实际位置与初始估计位置均在期望位置的基础上叠加了服从均匀分布的随机扰动。此外，为了模拟真实传感器测量的噪声特性，在测距信息中引入了均值为0、方差为0.03的高斯白噪声。

\begin{table}[H]
	\centering
	\caption{无人机集群系统刚性编队仿真基础参数表}
	\label{无人机集群系统刚性编队仿真基础参数表}
	\renewcommand{\arraystretch}{1.2}
	\begin{tabular}{lcc}
		\toprule
		\textbf{参数名称} & \textbf{符号} & \textbf{数值/配置} \\
		\midrule
		无人机数量 & $N$ & 4 \\
		通信与量测边数 & $M$ & 6 \\
		避碰安全距离下界 & $\rho_{e_1}$ & 0.5 m \\
		通信连通距离上界 & $\rho_{e_2}$ & 2.0 m \\
		期望机间距离 & $\sqrt{\beta_{ij}^*}$ & $\sqrt{2} \approx 1.414$ m \\
		仿真总时长 & $t_{\text{end}}$ & 5.0 s \\
		仿真步长频率 & $t_{\text{frequence}}$ & 1000 Hz \\
		机1量测配置 & $\tilde{C}_1$ & $\text{diag}(0, 0, 1)$ \\
		机2量测配置 & $\tilde{C}_2$ & $\text{diag}(0, 0, 0)$ \\
		机3量测配置 & $\tilde{C}_3$ & $\text{diag}(1, 1, 0)$ \\
		机4量测配置 & $\tilde{C}_4$ & $\text{diag}(1, 1, 1)$ \\
		增益参数 & $k_c,k_o$ & $k_c=5,k_o=10$ \\
		\bottomrule
	\end{tabular}
\end{table}

在验证刚性编队局部收敛之后，进一步给出四个典型编队变换场景以验证方法的实用性，其与刚性编队局部收敛实例由相同的参数配置与期望编队位置，仅是初始位置不同，具体设置见表 。

\begin{table}[htbp]
	\centering
	\caption{无人机典型编队变换场景初始状态设置}
	\label{tab:formation_transformation}
	\renewcommand{\arraystretch}{1.5}
	\begin{tabular}{c c m{8cm}}
		\toprule
		\textbf{场景编号} & \textbf{变换类型} & \textbf{初始位置参数设置} $\mathbf{p}_i(0)$ \\
		\midrule
		1 & 平移 (Translation) & $\mathbf{p}_i(0) = \mathbf{p}_i^* - [5, 5, 5]^T$，即初始位置由期望构型整体向负方向平移 5 m。 \\
		2 & 旋转 (Rotation) & $\mathbf{p}_i(0) = \mathbf{R}(\alpha, \beta, \gamma)\mathbf{p}_i^*$，其中欧拉角设置为 $\alpha = \frac{\pi}{2}, \beta = \frac{\pi}{2}, \gamma = \frac{\pi}{2}$。 \\
		3 & 扩张 (Expansion) & $\mathbf{p}_i(0) = 0.5 \mathbf{p}_i^*$，即初始编队尺度为期望构型的一半，随后向外扩张。 \\
		4 & 收缩 (Shrinkage) & $\mathbf{p}_i(0) = 1.5 \mathbf{p}_i^*$，即初始编队尺度为期望构型的 1.5 倍，随后向内收缩。 \\
		\bottomrule
	\end{tabular}
\end{table}


\subsection{编队轨迹与机间距离收敛性验证}

\subsection{状态估计与控制误差演化分析}

\section{Gazebo物理仿真平台验证}


















\section{字体效果表格}

% 公式上下不要空行，置于同一个段落下即可，否则上下距离会出现高度不一致的问题
\begin{equation}
	LRI=1\ ∕\ \sqrt{1+{\left(\frac{{\mu }_{R}}{{\mu }_{s}}\right)}^{2}{\left(\frac{{\delta }_{R}}{{\delta }_{s}}\right)}^{2}}
\end{equation}

% 列：Regular、Italic、Bold、Bold Italic
% 行：宋体、黑体、楷体、Serif、Sans Serif、Typewriter、Math

\begin{table}[htb]
	\centering
	\caption{字体效果表格}
	\begin{tabular}{@{}lllll@{}}
		\toprule
		& Regular & Bold & Italic & Bold Italic \\ \midrule
		宋体       & 宋体      & \colorbox{orange}{\textbf{宋体粗体}} & \textit{楷体}     &   \colorbox{gray}{\textbf{\textit{楷书粗斜体}}}  \\
		黑体         & {\heiti{}黑体}      & \textbf{\heiti{}黑体粗体} &   \textit{\heiti{}黑体斜体}     & \colorbox{gray}{\textit{\textbf{\heiti{}黑体粗斜体}}}   \\
		楷体         & {\kaishu{}楷书}      & \textbf{\kaishu{}楷书粗体} & \textit{\kaishu{}斜体楷体} &  \colorbox{gray}{\textbf{\textit{\kaishu{}楷书粗斜体}}}    \\
		Serif(Roman/Normal)      &    Regular    &  \textbf{Bold}  &    \textit{Italic}    &     \textbf{\textit{Bold Italic}}    \\
		Sans Serif &  \textsf{Regular}       &  \textbf{\textsf{Bold}}    &  \textit{\textsf{Bold}}   &    \textbf{\textit{\textsf{Bold}}}   \\
		% 有些字体缺少特定变体，LaTeX会抛出警告，同时尽量替换为相近变体。
		% 若编译结果能接受，可忽略之。
		% 例如此处Typewriter代表的Latin Modern Sans Typewriter（lmtt）字体没有bold extended italic（bx/it）变体，
		% 所以`\textbf{\textit{\texttt{…}}}`会被替换为bold slant（b/sl）变体，并引发如下警告。（其中TU指字体编码）
		% Font shape `TU/lmtt/bx/it' in size <…> not available. Font shape `TU/lmtt/b/sl' tried instead.
		Typewriter &  \texttt{Regular}       &  \textbf{\texttt{Bold}}    &  \textit{\texttt{Bold}}   &    \textbf{\textit{\texttt{Bold}}}   \\
		Math       &   $\mathnormal{Regular} \mathrm{Roman}$  & $\mathbf{Bold}$   &    $\mathit{Italic}$    &  $\mathbf{\mathit{Bold Italic}}$    \\ \bottomrule
	\end{tabular}
\end{table}

\begin{itemize}[nosep]
	\item \colorbox{orange}{宋体粗体}在 Windows 下会成为黑体。这是因为 Windows 的中易宋体没有粗体字重而进行的妥协。
	如果想要获得宋体粗体的样式，请在配置中开启伪粗体选项。
	\item \colorbox{gray}{粗斜体}的效果是因操作系统字体而定的，中文写作中不会使用这种字形，可以忽略。
\end{itemize}

\textit{有关公式与上下文间距的一些注意事项：请保证源码中的公式的环境（如}
\\ \verb|\begin{equation}|
	\textit{）与上一段落不要有空行。否则，公式和上文段落之间会有额外的空白。}
	
	
	\section{常见问题和疑难解答}
	
	如果您遇到\href{https://bithesis.bitnp.net/faq/char-missing.html}{生僻字无法显示}、
	\href{https://bithesis.bitnp.net/faq/enumitem-nosep.html}{列表项间距过大}、
	\href{https://bithesis.bitnp.net/faq/longtable.html}{三线表需要跨页}等问题，
	请参考\textcolor{magenta}{\href{https://bithesis.bitnp.net/faq/}{在线文档的「疑难杂症」部分}}。
	
	\section{代码片段}
	
	\begin{lstlisting}[language=Python, caption={Python Code}, label={lst:pythonfile}]
		import numpy as np
		
		def incmatrix(genl1,genl2):
		m = len(genl1)
		n = len(genl2)
		M = None #to become the incidence matrix
		VT = np.zeros((n*m,1), int)  #dummy variable
		
		#compute the bitwise xor matrix
		M1 = bitxormatrix(genl1)
		M2 = np.triu(bitxormatrix(genl2),1)
		
		for i in range(m-1):
		for j in range(i+1, m):
		[r,c] = np.where(M2 == M1[i,j])
		for k in range(len(r)):
		VT[(i)*n + r[k]] = 1;
		VT[(i)*n + c[k]] = 1;
		VT[(j)*n + r[k]] = 1;
		VT[(j)*n + c[k]] = 1;
		
		if M is None:
		M = np.copy(VT)
		else:
		M = np.concatenate((M, VT), 1)
		
		VT = np.zeros((n*m,1), int)
		
		return M
	\end{lstlisting}

